\documentclass[main.tex]{subfiles}


\begin{document}

%from pagenumber to up
% \phantomsection 
% \hypertarget{toc}{}

 

 

 
   

\section{Эквипотенциальные поверхности}


Характеризовать распределение поля в пространстве можно с помощью так называемых эквипотенциальных поверхностей.  

\textbf{Эквипотенциальная поверхность}~--- это поверхность в пространстве, все точки которой имеют \emph{одинаковый потенциал} ($\varphi=\mathrm{const}$). 

На рис.\,\ref{fig:p102} изображены эквипотенциальные поверхности поля для двух случаев: точечный заряд и заряженная пластина.

    \begin{figure}[H]
    \centering

    \begin{tikzpicture}[font=\small,            line cap=round,                             >={Stealth[inset=0pt,round,length=3mm, angle'=20]},vec/.style={>={Stealth[inset=0pt,round,length=3.5mm, angle'=20]}}]


\begin{scope}[yshift=1.5cm]

%pot line
\foreach \i in{1,...,3}
\draw[gray] (0,0) circle ({.5*\i});

\shade[ball color=red] (0,0) circle (.15);

%\phi
% \node[below] at (0,.5) {$\varphi$};

\end{scope}

%%%%%%%%%%%%%%%%%%%%%%

\begin{scope}[xshift=6cm]

%pot line
\foreach \i in{1,...,3}
\draw[gray,xshift={.5*\i cm}] (0,0) --++ (0,3);

\foreach \i in{-1,...,-3}
\draw[gray,xshift={.5*\i cm}] (0,0) --++ (0,3);


%plane
\draw[thick,red] (0,0) --++ (0,3);
    
\end{scope}


    \end{tikzpicture}
\caption{Эквипотенциальные поверхности поля заряженного тела}
\label{fig:p102}
\end{figure}

Эквипотенциальными поверхностями поля точечного заряда являются всевозможные сферы с центром в заряде (серые линии на рис.\,\ref{fig:p102}, слева). Поле же вблизи равномерно заряженной пластины описывается эквипотенциальными поверхностями, параллельными пластине (серые линии на рис.\,\ref{fig:p102}, справа).

\emph{Эквипотенциальные поверхности всегда перпендикулярны линиям поля.}\linebreak Этот факт иллюстрируется рисунком~\ref{fig:p103}. 

    \begin{figure}[H]
    \centering

    \begin{tikzpicture}[font=\small,            line cap=round,                             >={Stealth[inset=0pt,round,length=3mm, angle'=20]},vec/.style={>={Stealth[inset=0pt,round,length=3.5mm, angle'=20]}}]


\begin{scope}[yshift=1.5cm]

%line
\foreach \i in{0,...,11}
\draw[->,Green] (0,0) --++ ({30*\i}:2);

%pot line
\foreach \i in{1,...,3}
\draw[gray] (0,0) circle ({.5*\i});

\shade[ball color=red] (0,0) circle (.15);

%\phi
% \node[below] at (0,.5) {$\varphi$};

\end{scope}

%%%%%%%%%%%%%%%%%%%%%%

\begin{scope}[xshift=6cm]

%line
\foreach \i in{0,...,5}
\draw[->,Green] (0,{.25+\i*.5}) --++ ({180}:2);

\foreach \i in{0,...,5}
\draw[->,Green] (0,{.25+\i*.5}) --++ ({0}:2);


%pot line
\foreach \i in{1,...,3}
\draw[gray,xshift={.5*\i cm}] (0,0) --++ (0,3);

\foreach \i in{-1,...,-3}
\draw[gray,xshift={.5*\i cm}] (0,0) --++ (0,3);


%plane
\draw[thick,red] (0,0) --++ (0,3);
    
\end{scope}


    \end{tikzpicture}
\caption{Перпендикулярность эквипотенциальных поверхностей и линий поля}
\label{fig:p103}
\end{figure}


\noindent\textbf{Задача.} На рис.\,\ref{fig:p104} показаны линии однородного электрического поля и две эквипотенциальные поверхности (А и~Б). Какая поверхность, А или~Б, характеризуется большим значением потенциала?

    \begin{figure}[H]
    \centering

    \begin{tikzpicture}[font=\small,            line cap=round,                             >={Stealth[inset=0pt,round,length=3mm, angle'=20]},vec/.style={>={Stealth[inset=0pt,round,length=3.5mm, angle'=20]}}]


\begin{scope}[xshift=6cm]

%line
\foreach \i in{0,...,2}
\draw[Green,decoration={markings,mark=at position {1/2+1.5/20} with {\arrow{>}}},postaction={decorate}] (0,{.25+\i*.5}) --++ ({0}:2);


%pot line
\draw[gray,xshift={.5*1 cm}] (0,0) --++ (0,1.5) node[above,black] {А};
\draw[gray,xshift={.5*3 cm}] (0,0) --++ (0,1.5) node[above,black] {Б};



    
\end{scope}


    \end{tikzpicture}
\caption{К задаче}
\label{fig:p104}
\end{figure}

\noindent\emph{Решение}. Как известно, в направлении линии поля потенциал поля убывает. Значит, потенциал поверхности~А больше потенциала поверхности~Б. 

























 
\end{document}